% Paper B — Section on institutional review cadence (survey).
% House style: literal author-year prose citations, no citation macros.

\section{The cadence institutions actually use}
\label{sec:cadence}

The decision clock is not an abstraction imposed on governance from outside. It is a live,
varying, and \emph{revised} institutional choice. This section surveys the review cycles
that real management bodies operate on, so that the horizons derived earlier can be set
against the intervals they have to sit inside.

\subsection{The survey}

\begin{itemize}

\item \textbf{International Whaling Commission, aboriginal subsistence whaling.}
Catch quotas --- known as \emph{strike limits} --- are set in \textbf{six-year blocks},
established by Strike Limit Algorithms tested by simulation over a 100-year period. The
current blocks were established at the Commission meeting of September 2018, renewed in
2024, and fall due for review in 2030. \emph{The cadence itself is the notable fact}: until
the multi-year process was adopted in 2018, strike limits required a vote at a Commission
meeting \textbf{every two years}. The 2018 reform replaced that with an automatic six-year
rollover, conditional on three tests --- that the number of whales needed has not
increased, that the Scientific Committee advises the limits will not harm the stocks, and
that the Commission determines the governments concerned have complied with the agreed
timeline and that the information supplied represents a status quo continuation. The
rollover was used successfully for the first time in 2024.

\item \textbf{NOAA Fisheries, United States.} Management track assessments are held
\textbf{twice a year}, usually in June and September, but each stock has its own
established update cycle: some are updated \textbf{every year} and others have more than a
year between updates. Most species in the Mid-Atlantic and New England regions sit on a
\textbf{two-year} cycle; in Hawai`i the cycle for most island species is \textbf{three
years}. The more comprehensive research or benchmark tracks are performed less frequently
--- approximately \textbf{every five years} for a regionally important stock such as red
snapper in the southeast, while \textbf{a decade} may pass before other stocks are
reassessed. NOAA provides scientific information for roughly 500 stocks and has capacity to
assess about 200 each year. Marine mammal stocks are reviewed \textbf{annually} where
strategic and \textbf{every three years} where not.

\item \textbf{ICES (International Council for the Exploration of the Sea).} The advice is
based on \textbf{annual} fish stock assessments, but the \emph{methodology} behind those
assessments is evaluated in a benchmark workshop \textbf{every three to five years}, with
roughly twenty stocks benchmarked per year and ACOM agreeing the benchmark programme a year
in advance. Preparation alone takes five to seven months. The two clocks are therefore
different and both matter: the advice moves yearly, the assessment machinery that produces
it moves on a three-to-five-year cycle.

\item \textbf{Marine Stewardship Council certification.} Full certification is granted for
\textbf{five years}, subject to \textbf{annual} surveillance audits, with a full
re-assessment at the end of the fifth year. Certification may be suspended or withdrawn
within the cycle if a surveillance audit finds a material change. Chain-of-custody
certification runs on a separate, shorter clock: three years, with surveillance audits at
twelve- or eighteen-month intervals.

\end{itemize}

\subsection{The cadence against the horizon}

The point of setting these side by side is not to certify any of them. It is to show that
the interval a real institution operates on is of the same order as --- and in several
cases at or beyond --- the horizon a certificate of the kind derived in Section
\ref{sec:cod} can carry.

\begin{center}
\begin{tabular}{lll}
\hline
Institution & Cadence & Against a 6--7 yr certified horizon \\
\hline
NOAA management track (most stocks) & 1--2 yr & well inside \\
ICES advice & 1 yr & well inside \\
NOAA research track, important stocks & $\sim$5 yr & inside, with little margin \\
ICES benchmark & 3--5 yr & inside, with little margin \\
MSC certification & 5 yr, annual audit & inside; audit is the fast clock \\
IWC strike limits (since 2018) & 6 yr & \textbf{at the boundary} \\
NOAA research track, other stocks & up to a decade & \textbf{beyond} \\
\hline
\end{tabular}
\end{center}

Two observations follow, and both are stated with the scope they actually have.

The first is that the fast clocks and the slow clocks in these institutions are doing
different jobs, and the distinction matters for where the decision clock sits. ICES advice
moves annually while its benchmark moves every three to five years; MSC certification runs
for five years while its surveillance runs annually. An institution is rarely
single-clocked. What has to be compared against a certification horizon is the clock that
governs \emph{what is held fixed} --- the interval over which a control is applied without
revision --- and not the clock that governs how often it is looked at.

The second is the one that carries the design lesson, and it comes from the IWC. That body
did not merely operate on a cadence; it \textbf{changed} one, deliberately, from two years
to six, and it did so by making the rollover automatic subject to stated tests. The
instrument of reform was not a new scientific assessment but a change in the rules by which
the existing assessment is renewed. That is the clearest available demonstration that the
decision clock is a design variable in the sense this paper has been using the term: an
institution treated it as one, moved it, and said why.

\subsection{Scope, stated plainly}

This table is a survey of cadences, not a set of certifications. The certified horizon of
six to seven years derived in Section \ref{sec:cod} is a property of one stock --- Northern
cod, under a stated fit, a stated protocol and a stated admissible set for \(K\) --- and it
does not transfer to any other stock by inspection. What the comparison shows is that real
institutional cadences are of the same order of magnitude as that horizon, so that the
question ``does this institution's clock sit inside the horizon its certificate can
carry?'' is a live one and not a hypothetical. Answering it for any particular stock
requires the certificate for that stock. It cannot be read off this table.

\subsection{Gaps in this survey}

The survey covers four bodies and is not exhaustive. Not covered, and needed before this
section is complete: the tuna RFMOs (ICCAT, IATTC, WCPFC, CCSBT) and their
commission-meeting and assessment cycles; the European Union's Common Fisheries Policy,
where total allowable catches are set annually while multiannual plans run on longer
cycles; national agencies other than NOAA, notably Fisheries and Oceans Canada; and the
Southern Hemisphere harvest-strategy frameworks of Australia and New Zealand. These are
listed as open rather than assumed to resemble the four surveyed.
